% GENERATED FILE. Edit canonical lessons, then run npm run generate:books.
% canonical-source-hash: fnv1a64:9bff0221cce4d5d6
% canonical-lessons: BN-C395-phlait, BN-C395-ranoye, BN-C395-kastoms, BN-C395-bidesh, BN-C395-pari-deoya

\chapter{Flights and Runways}
\label{ch:bn-flights-and-runways}
\hlchaptermodality{395}

\begin{chapteropening}
\textbf{By the end of this chapter:} \emph{I can talk about a flight, the runway and customs on the way abroad.}

Complete the last lesson of chapter 395: I can talk about a flight, the runway and customs on the way abroad.
\end{chapteropening}

\section[\texorpdfstring{\bn{ফ্লাইট} (phlāiṭ)}{ফ্লাইট (phlāiṭ)}]{\bn{ফ্লাইট} (phlāiṭ) — a flight}

\label{lesson:BN-C395-phlait}



\begin{quote}
Before the new one: say the Bengali for a steering wheel, then the Bengali for to rub.
\end{quote}

\subsection*{You'll want to know: \bn{ফ্লাইট}}
\textbf{\bn{ফ্লাইট}} — \emph{phlāiṭ} — \textquotedblleft{}a flight\textquotedblright{}.

\begin{quote}
\bn{আমাদের} \bn{ফ্লাইট} \bn{রাত} \bn{দশটায়} \bn{ছাড়বে।}
\end{quote}

\begin{quote}
\emph{āmāder phlāiṭ rāt dôshṭāy chhāṛbe.}
\end{quote}

\begin{quote}
Our flight leaves at ten at night.
\end{quote}

\begin{cousinweb}
\textbf{\bn{ফ্লাইট}} (\emph{phlāiṭ}) is borrowed from English \emph{flight}, with \bn{ফ} for f and the hard \bn{ট} for t.
\end{cousinweb}

\begin{grammarlens}[title={What you've built}]
One more word for getting around.
\end{grammarlens}

\subsection*{Your turn}
\begin{itemize}
\raggedright
  \item \emph{Say it:} \emph{phlāiṭ}
  \item \emph{Say it:} \emph{phlāiṭ}, once more
  \item \emph{Say it:} \emph{āmāder phlāiṭ rāt dôshṭāy chhāṛbe.} — Our flight leaves at ten at night.
  \item \emph{Recall:} say the Bengali for a steering wheel, then the Bengali for to rub, then say \emph{phlāiṭ} again
\end{itemize}

\begin{tcolorbox}[breakable,colback=teal!4,colframe=teal!35!black,title={Before you move on}]
What does \bn{ফ্লাইট} mean? (\textquotedblleft{}A flight\textquotedblright{}.) Say it once more.
\end{tcolorbox}
\section[\texorpdfstring{\bn{রানওয়ে} (rānoye)}{রানওয়ে (rānoye)}]{\bn{রানওয়ে} (rānoye) — a runway}

\label{lesson:BN-C395-ranoye}



\begin{quote}
Before the new one: say the Bengali for to rub, then the Bengali for a flight.
\end{quote}

\subsection*{You'll want to know: \bn{রানওয়ে}}
\textbf{\bn{রানওয়ে}} — \emph{rānoye} — \textquotedblleft{}a runway\textquotedblright{}.

\begin{quote}
\bn{বিমানটা} \bn{ধীরে} \bn{রানওয়েতে} \bn{নামল।}
\end{quote}

\begin{quote}
\emph{bimānṭā dhīre rānoyete nāmlô.}
\end{quote}

\begin{quote}
The plane came down slowly onto the runway.
\end{quote}

\begin{cousinweb}
\textbf{\bn{রানওয়ে}} (\emph{rānoye}) is borrowed from English \emph{runway}; the w is written with \bn{ও} and \bn{য়}.
\end{cousinweb}

\begin{grammarlens}[title={What you've built}]
One more word for getting around.
\end{grammarlens}

\subsection*{Your turn}
\begin{itemize}
\raggedright
  \item \emph{Say it:} \emph{rānoye}
  \item \emph{Say it:} \emph{rānoye}, once more
  \item \emph{Say it:} \emph{bimānṭā dhīre rānoyete nāmlô.} — The plane came down slowly onto the runway.
  \item \emph{Recall:} say the Bengali for to rub, then the Bengali for a flight, then say \emph{rānoye} again
\end{itemize}

\begin{tcolorbox}[breakable,colback=teal!4,colframe=teal!35!black,title={Before you move on}]
What does \bn{রানওয়ে} mean? (\textquotedblleft{}A runway\textquotedblright{}.) Say it once more.
\end{tcolorbox}
\section[\texorpdfstring{\bn{কাস্টমস} (kāsṭôms)}{কাস্টমস (kāsṭôms)}]{\bn{কাস্টমস} (kāsṭôms) — customs (at a border or airport)}

\label{lesson:BN-C395-kastoms}



\begin{quote}
Before the new one: say the Bengali for a flight, then the Bengali for a runway.
\end{quote}

\subsection*{You'll want to know: \bn{কাস্টমস}}
\textbf{\bn{কাস্টমস}} — \emph{kāsṭôms} — \textquotedblleft{}customs (at a border or airport)\textquotedblright{}.

\begin{quote}
\bn{কাস্টমসে} \bn{আমাদের} \bn{ব্যাগ} \bn{খুলতে} \bn{হল।}
\end{quote}

\begin{quote}
\emph{kāsṭômse āmāder byāg khulte hôlô.}
\end{quote}

\begin{quote}
At customs we had to open our bags.
\end{quote}

-\bn{তে} \bn{হল} is the past of -\bn{তে} \bn{হবে}: \bn{খুলতে} \bn{হল}, we had to open.

\begin{cousinweb}
\textbf{\bn{কাস্টমস}} (\emph{kāsṭôms}) is borrowed from English \emph{customs}, the office at a border that checks goods and collects duty.
\end{cousinweb}

\begin{grammarlens}[title={What you've built}]
One more word for getting around.
\end{grammarlens}

\subsection*{Your turn}
\begin{itemize}
\raggedright
  \item \emph{Say it:} \emph{kāsṭôms}
  \item \emph{Say it:} \emph{kāsṭôms}, once more
  \item \emph{Say it:} \emph{kāsṭômse āmāder byāg khulte hôlô.} — At customs we had to open our bags.
  \item \emph{Recall:} say the Bengali for a flight, then the Bengali for a runway, then say \emph{kāsṭôms} again
\end{itemize}

\begin{tcolorbox}[breakable,colback=teal!4,colframe=teal!35!black,title={Before you move on}]
What does \bn{কাস্টমস} mean? (\textquotedblleft{}Customs (at a border or airport)\textquotedblright{}.) Say it once more.
\end{tcolorbox}
\section[\texorpdfstring{\bn{বিদেশ} (bidesh)}{বিদেশ (bidesh)}]{\bn{বিদেশ} (bidesh) — abroad; a foreign country}

\label{lesson:BN-C395-bidesh}



\begin{quote}
Before the new one: say the Bengali for a runway, then the Bengali for customs.
\end{quote}

\subsection*{You'll want to know: \bn{বিদেশ}}
\textbf{\bn{বিদেশ}} — \emph{bidesh} — \textquotedblleft{}abroad; a foreign country\textquotedblright{}.

\begin{quote}
\bn{মীরার} \bn{ভাই} \bn{বিদেশে} \bn{থাকে।}
\end{quote}

\begin{quote}
\emph{mīrār bhāi bideshe thāke.}
\end{quote}

\begin{quote}
Mira's brother lives abroad.
\end{quote}

\bn{বিদেশে}, abroad (being there); \bn{বিদেশে} \bn{যাওয়া}, to go abroad. \bn{বিদেশি}, foreign, is built on it.

\begin{cousinweb}
\textbf{\bn{বিদেশ}} (\emph{bidesh}) is Sanskrit \emph{videśa}, \textquotedblleft{}another country, abroad\textquotedblright{} (Monier-Williams), from \emph{vi-}, apart, and \emph{deśa}, a country: the land away from home.
\end{cousinweb}

\begin{grammarlens}[title={What you've built}]
One more word for getting around.
\end{grammarlens}

\subsection*{Your turn}
\begin{itemize}
\raggedright
  \item \emph{Say it:} \emph{bidesh}
  \item \emph{Say it:} \emph{bidesh}, once more
  \item \emph{Say it:} \emph{mīrār bhāi bideshe thāke.} — Mira's brother lives abroad.
  \item \emph{Recall:} say the Bengali for a runway, then the Bengali for customs, then say \emph{bidesh} again
\end{itemize}

\begin{tcolorbox}[breakable,colback=teal!4,colframe=teal!35!black,title={Before you move on}]
What does \bn{বিদেশ} mean? (\textquotedblleft{}Abroad; a foreign country\textquotedblright{}.) Say it once more.
\end{tcolorbox}
\section[\texorpdfstring{\bn{পাড়ি} \bn{দেওয়া} (pāṛi deoyā)}{পাড়ি দেওয়া (pāṛi deoyā)}]{\bn{পাড়ি} \bn{দেওয়া} (pāṛi deoyā) — to make a crossing (of sea or river); to set out on a long way}

\label{lesson:BN-C395-pari-deoya}



\begin{quote}
Before the new one: say the Bengali for customs, then the Bengali for abroad.
\end{quote}

\subsection*{You'll want to know: \bn{পাড়ি} \bn{দেওয়া}}
\textbf{\bn{পাড়ি} \bn{দেওয়া}} — \emph{pāṛi deoyā} — \textquotedblleft{}to make a crossing (of sea or river); to set out on a long way\textquotedblright{}.

\begin{quote}
\bn{জাহাজটা} \bn{সাত} \bn{দিনে} \bn{সমুদ্র} \bn{পাড়ি} \bn{দিল।}
\end{quote}

\begin{quote}
\emph{jāhājṭā sāt dine sômudrô pāṛi dilô.}
\end{quote}

\begin{quote}
The ship crossed the sea in seven days.
\end{quote}

Past (he/she …): \bn{সেই} \bn{জাহাজ} \bn{সমুদ্র} \bn{পাড়ি} \bn{দিল}, \emph{sei jāhāj sômudrô pāṛi dilô}.

A word with distance in it: you \bn{পাড়ি} \bn{দাও} across an ocean, or to a far city.

\begin{cousinweb}
\textbf{\bn{পাড়ি} \bn{দেওয়া}} (\emph{pāṛi deoyā}) joins \bn{পাড়ি}, \textquotedblleft{}a crossing, a long passage\textquotedblright{}, to \bn{দেওয়া}, to give; the origin of \bn{পাড়ি} is not settled.
\end{cousinweb}

\begin{grammarlens}[title={What you've built}]
One more word for getting around.
\end{grammarlens}

\subsection*{Your turn}
\begin{itemize}
\raggedright
  \item \emph{Say it:} \emph{pāṛi deoyā}
  \item \emph{Say it:} \emph{pāṛi deoyā}, once more
  \item \emph{Say it:} \emph{jāhājṭā sāt dine sômudrô pāṛi dilô.} — The ship crossed the sea in seven days.
  \item \emph{Recall:} say the Bengali for customs, then the Bengali for abroad, then say \emph{pāṛi deoyā} again
\end{itemize}

\begin{tcolorbox}[breakable,colback=teal!4,colframe=teal!35!black,title={Before you move on}]
What does \bn{পাড়ি} \bn{দেওয়া} mean? (\textquotedblleft{}To make a crossing (of sea or river); to set out on a long way\textquotedblright{}.) Say it once more.
\end{tcolorbox}
